
\begin{center}
{\LARGE\bfseries\color{azulPDMI}
Quantos triângulos se escondem num polígono?}\\[2pt]
{\small Soma das amplitudes dos ângulos internos
-- tarefa de exploração}
\end{center}

\vspace{2pt}

\noindent
\textbf{Nome:} \rule{8.1cm}{0.4pt}
\hfill
\textbf{N.º:} \rule{1.3cm}{0.4pt}

\begin{tcolorbox}[
colback=azulClaro,
colframe=azulPDMI!45,
boxrule=0.5pt,
arc=2mm,
left=7pt,
right=7pt,
top=5pt,
bottom=5pt
]

\textbf{O desafio}\par

Sabemos que a soma das amplitudes dos
ângulos internos de um triângulo é
$180^\circ$.

Conseguiremos descobrir a soma dos
ângulos internos de um
\textbf{polígono convexo sem medir
nenhum ângulo}?

Regista desenhos, cálculos e explicações
que permitam compreender o teu raciocínio.

\end{tcolorbox}

% ========================================================
% PARTE I
% ========================================================

\section*{Parte I -- À procura de uma estratégia}

\textbf{1.1.}
Observa o quadrilátero e o pentágono.
Procura uma estratégia para obter a soma
das amplitudes dos ângulos internos de
cada um, \textbf{sem utilizar o transferidor}.

Podes desenhar sobre as figuras.
Explica como pensaste.

\vspace{2pt}

\begin{center}

\begin{minipage}{.45\linewidth}
\centering

\begin{tikzpicture}[scale=.65,line join=round]

\coordinate (A) at (.3,.2);
\coordinate (B) at (4.3,.1);
\coordinate (C) at (4.7,2.6);
\coordinate (D) at (.5,2.9);

\draw[
thick,
draw=azulPDMI,
fill=azulClaro
]
(A)--(B)--(C)--(D)--cycle;

\node[below left] at (A) {$A$};
\node[below right] at (B) {$B$};
\node[above right] at (C) {$C$};
\node[above left] at (D) {$D$};

\end{tikzpicture}

\vspace{2pt}

{\small Quadrilátero $ABCD$}

\end{minipage}
\hfill
\begin{minipage}{.45\linewidth}
\centering

\begin{tikzpicture}[scale=.60,line join=round]

\coordinate (A) at (.4,.2);
\coordinate (B) at (3.8,.1);
\coordinate (C) at (5,2.2);
\coordinate (D) at (2.8,3.4);
\coordinate (E) at (.1,2);

\draw[
thick,
draw=azulPDMI,
fill=azulClaro
]
(A)--(B)--(C)--(D)--(E)--cycle;

\node[below left] at (A) {$A$};
\node[below right] at (B) {$B$};
\node[right] at (C) {$C$};
\node[above] at (D) {$D$};
\node[left] at (E) {$E$};

\end{tikzpicture}

\vspace{2pt}

{\small Pentágono $ABCDE$}

\end{minipage}

\end{center}

\vspace{-5pt}

\linhas{2}

\textbf{1.2.}
Compara a tua ideia com a de um colega.
Chegaram ao mesmo resultado?
Identificaram estratégias diferentes?
Partilhem uma explicação.

\linhas{1}

\begin{tcolorbox}[
colback=azulClaro,
colframe=azulPDMI!45,
boxrule=.5pt,
arc=2mm,
left=7pt,
right=7pt,
top=5pt,
bottom=5pt
]

\textbf{Uma estratégia para experimentar}

Uma \textbf{diagonal} une dois vértices
não consecutivos de um polígono.

Mesmo que tenhas encontrado outra
estratégia, experimenta agora:

\begin{enumerate}[
label=\alph*),
leftmargin=1.7em,
nosep,
topsep=2pt
]

\item
Traça uma diagonal no quadrilátero.
Em que figuras ficou dividido?

\item
No pentágono, escolhe um vértice
e traça todas as diagonais que
partem dele.
Quantos triângulos obténs?

\item
Usa a soma dos ângulos internos
de um triângulo para calcular a
soma dos ângulos internos de
cada polígono.

\end{enumerate}

\end{tcolorbox}

\textbf{1.3.}
Porque é que podemos somar os ângulos
dos triângulos para obter a soma dos
ângulos internos do polígono?

Na tua explicação, considera se há
sobreposições ou regiões por preencher.

\linhas{2}


\section*{Parte II -- Da exploração à regra geral}

\textbf{2.1.}
Completa a tabela.

Em cada polígono, desenha as diagonais
\textbf{a partir de um único vértice},
de modo a obter triângulos sem
sobreposições e sem espaços por preencher.

\begin{center}

\renewcommand{\arraystretch}{1.12}

\begin{tabularx}{\linewidth}{|
>{\centering\arraybackslash}p{3.5cm}|
>{\centering\arraybackslash}p{1.6cm}|
>{\centering\arraybackslash}p{2.2cm}|
>{\centering\arraybackslash}X|}

\hline

\rowcolor{azulClaro}

\textbf{Polígono} &
\textbf{N.º de lados} &
\textbf{N.º de triângulos} &
\textbf{Soma dos ângulos internos}

\\ \hline

% TRIANGULO

\parbox[c][1.92cm][c]{3.5cm}{
\centering

\begin{tikzpicture}[
scale=.72,
baseline=(current bounding box.center)
]

\draw[
thick,
azulPDMI,
fill=azulClaro
]
(0,0)--(2.8,0)--(1.4,1.7)--cycle;

\end{tikzpicture}

\vspace{3pt}

{\scriptsize Triângulo}
}

& 3 & 1 & $180^\circ$

\\ \hline

% QUADRILATERO

\parbox[c][1.92cm][c]{3.5cm}{
\centering

\begin{tikzpicture}[
scale=.69,
baseline=(current bounding box.center)
]

\draw[
thick,
azulPDMI,
fill=azulClaro
]
(0,0)--(2.8,.1)--(2.7,1.8)--
(.2,1.7)--cycle;

\end{tikzpicture}

\vspace{3pt}

{\scriptsize Quadrilátero}
}

& 4 & &

\\ \hline

% PENTAGONO

\parbox[c][1.92cm][c]{3.5cm}{
\centering

\begin{tikzpicture}[
scale=.69,
baseline=(current bounding box.center)
]

\draw[
thick,
azulPDMI,
fill=azulClaro
]
(0,.6)--(1,.05)--(2.7,.2)--
(3,1.6)--(1.2,2)--cycle;

\end{tikzpicture}

\vspace{3pt}

{\scriptsize Pentágono}
}

& 5 & &

\\ \hline

% HEXAGONO

\parbox[c][1.92cm][c]{3.5cm}{
\centering

\begin{tikzpicture}[
scale=.69,
baseline=(current bounding box.center)
]

\draw[
thick,
azulPDMI,
fill=azulClaro
]
(.2,.4)--(1,.05)--(2.4,.2)--
(3,1.1)--(2.1,2)--(.3,1.7)--cycle;

\end{tikzpicture}

\vspace{3pt}

{\scriptsize Hexágono}
}

& 6 & &

\\ \hline

\end{tabularx}

\end{center}

\renewcommand{\arraystretch}{1.55}

\textbf{2.2.}
Que relação observas entre o número
de lados do polígono e o número
de triângulos obtidos?

Descreve-a por palavras.

\linhas{1}

\textbf{2.3.}
Considera um polígono convexo
com $n$ lados.

Propõe uma expressão para a soma
$S$ das amplitudes dos seus
ângulos internos.

Explica o significado dos
fatores utilizados.

\linhas{2}
\newpage
\textbf{2.4.}
A tua regra funciona para
\emph{qualquer} polígono convexo?

Procura justificá-la sem te baseares
apenas nos exemplos da tabela.

Pensa no número de diagonais que
podes traçar a partir de um vértice
e no número de triângulos assim formado.

\linhas{2}

% ========================================================
% CONCLUSAO
% ========================================================

\begin{tcolorbox}[
colback=azulClaro,
colframe=azulPDMI!40,
boxrule=.5pt,
arc=2mm,
left=7pt,
right=7pt,
top=5pt,
bottom=5pt
]

\textbf{Para concluir}

Completa:

«Num polígono convexo com $n$ lados,
a soma das amplitudes dos ângulos
internos é

\[
S=\underline{\hspace{3.0cm}}
\]

porque
\underline{\hspace{4cm}}».

\end{tcolorbox}

\vspace{2pt}

{\footnotesize\itshape

\textbf{Se houver tempo:}

Sem desenhar, determina a soma das
amplitudes dos ângulos internos
de um polígono convexo com 12 lados.

Justifica.

}
